\documentclass[11pt]{article}
\usepackage{amsmath,amssymb,hyperref}
\usepackage{geometry}
\usepackage{bbm}
\geometry{margin=1in}

% Generated by wolfbook-docs-agent from the live notebook on 2026-07-15.
% Reflects notebook state at generation time -- regenerate after material
% changes rather than hand-editing this file. Content corresponds to the
% companion XXX_CG_L3.md, but is restructured into continuous, topic-organized
% prose rather than a per-symbol transliteration, and every Greek-letter
% variable name from the notebook (lambda, theta, kappa, chi, tau) is rendered
% as a proper LaTeX math macro rather than a literal Unicode glyph. This file
% compiles cleanly with plain pdflatex.

\title{XXX\_CG\_L3 --- Reference}
\author{}
\date{Generated by wolfbook-docs-agent from the live notebook on 2026-07-15}

\newcommand{\qth}{Q_\theta}

\begin{document}
\maketitle

\section*{Overview}

This is the second Clean notebook in the XXX subproject: the three-site ($L=3$) rational
$\mathfrak{gl}(2)$ (XXX) spin chain --- independent representations $[\lambda_1,0]$, $[\lambda_2,0]$,
$[\lambda_3,0]$ at each site, independent inhomogeneities $\theta_1,\theta_2,\theta_3$, and a shared
companion twist $G$ --- generalized directly from \texttt{XXX\_CG\_L2\_Clean.wb}, together with the
associated Baxter TQ-equation solver. Unlike L2, this 44-cell notebook was built \textbf{directly in
\texttt{Clean/}}, cell by cell, with no \texttt{Experiments/} merge or promotion step: it carries its
own complete validation suite, including the end-to-end TQ residual check that at L2 lives only in
the separate \texttt{Experiments/Baxter\_Solver\_L2.wb}. It follows the same two-section layout as
L2:

\begin{itemize}
\item \textbf{Section A} (cells 1--32) --- numeric parameters; GT-basis/generator machinery (ported
  verbatim from L2, itself from \texttt{Experiments/XXX\_su2.wb}, single-site and $L$-independent);
  three-site Lax/monodromy/transfer matrices; the fusion (Hirota) hierarchy; the SoV $B$-operator;
  and the SoV pseudovacuum/covector basis --- each followed by its own inline validation cell, same
  convention as L2.
\item \textbf{Section B} (cells 33--44) --- the Baxter/TQ solver: Baxter-equation ingredients,
  transfer-matrix coefficient extraction, the predicted weight-subspace dimension, the genuine
  transfer-matrix-eigenstate solve (\texttt{TauEigensystem}), accessors (\texttt{Psi}/$\tau$), the
  linear TQ solve (\texttt{QSolve}/\texttt{Q1}), and the end-to-end TQ-equation residual regression
  check.
\end{itemize}

Every construction in Section A with a direct $L=1$/$L=2$ counterpart is cross-referenced against
\texttt{paper/xxxCG.tex} \S2 below, generalized from the stated $L=1$ formulas exactly the way L2's
were; the paper does not state an $L=2$ or $L=3$ case explicitly, so these cross-references are
tensor-product generalizations, not literal transcriptions. Section B has no paper counterpart at
all --- \texttt{paper/xxxCG.tex} stops at the SoV pseudovacuum/covector construction and does not yet
state a TQ equation.

\paragraph{Why L3 needed real re-derivation, not just mechanical generalization.} At $L=2$ the
transfer matrix $t_{1,1}(u)$ is quadratic in $u$: one trivial leading coefficient plus exactly one
genuinely state-dependent coefficient, feeding a single free number $J$ per eigenstate. At $L=3$,
$t_{1,1}(u)$ is \textbf{cubic}: writing
\[
  t_{1,1}(u) = C_0 + C_1 u + C_2 u^2 + C_3 u^3,
\]
the leading coefficient $C_3\propto\mathbbm{1}$ is again trivial, and $C_2$ turns out to be a pure
Cartan/total-generator quantity (numerically confirmed in \S\ref{sec:asymptotics} below --- this is
the load-bearing cross-check the whole Baxter re-derivation depends on), but $C_1$ and $C_0$
genuinely mix different sites' raising and lowering operators and do \emph{not} reduce to closed
forms. They must instead be recorded directly, per eigenvector, as \emph{two} free numbers $(J_1,J_0)$
--- L2's single $J$ generalizes to a pair. This is what changes \texttt{TauEigensystem},
\texttt{BaxterMatrix}, and \texttt{QSolve} relative to their L2 counterparts, as detailed in
\S\ref{sec:baxter} below; everything else in Section A (the fusion hierarchy, \texttt{dPred},
\texttt{Q1}) generalizes mechanically, since the algebraic form of those constructions never
references $L$ explicitly.

\paragraph{A note on notation.} As in the notebook itself and its L2 predecessor, literal Greek
Unicode characters ($\lambda$, $\theta$, $\kappa$, $\chi$, $\tau$) are used as Wolfram symbol names;
throughout this document these are rendered as their corresponding math-mode macros. Only Wolfram
function calls with Latin/numeric arguments (e.g.\ \texttt{TauEigensystem[1,1,1]},
\texttt{QSolve[1,1,1][2,0]}) are quoted verbatim in \texttt{typewriter} font.

\section{The Gelfand--Tsetlin basis, single-site generators, and the companion twist}
\label{sec:gt}

The single-site building blocks --- Gelfand--Tsetlin (GT) basis enumeration (\texttt{ps[rep]},
memoized), irrep dimension (\texttt{dim[rep]}, memoized), the Cartan/raising/lowering generator
action on a GT pattern (\texttt{J0[k][pat]}, \texttt{Jp[k][pat]}, \texttt{Jm[k][pat]}), and the full
generator matrix assembly \texttt{JJ[k,l,rep]} --- are ported verbatim from \texttt{XXX\_CG\_L2\_Clean.wb}
Cell 3 (itself ported from \texttt{Experiments/XXX\_su2.wb}) and are independent of the chain length
$L$; every later construction is built on top of them. As at L2, this gives the expected
$\dim V_\lambda=\lambda+1$ for the two-row $\mathfrak{gl}(2)$ shapes used throughout, confirmed in
Cell 5: $\dim\{1,0\}=2$, $\dim\{2,0\}=3$, $\dim\{3,0\}=4$. All of this machinery is defined in Cell 4.

The wrapper $E_{ij}\in\mathrm{End}(V_\lambda)$ (routine \texttt{Ee[$\lambda$][i,j]}, Cell 6, simply
\texttt{JJ[i,j,\{$\lambda$,0\}]}) and the single-site Lax matrix (routine \texttt{L[$\lambda$][i,j][u]},
Cell 7),
\begin{equation}
  L^{[\lambda]}(u) = u\,\mathbbm{1} - h\sum_{i,j} E_{ji}\otimes e_{ij},
  \label{eq:Lax}
\end{equation}
match \texttt{paper/xxxCG.tex} \S2 exactly, again identically to L2; $u$ is left unshifted here, with
the inhomogeneity shift $u\mapsto u-\theta$ applied by the caller when a site is instantiated below.
Cell 8 confirms $\dim L^{[2]}_{11}(u)=3\times3$.

The companion-form twist matrix $G$ (Cell 9) is
\begin{equation}
  G = \begin{pmatrix} \chi_1 & -\chi_2 \\ 1 & 0 \end{pmatrix}, \qquad
  \chi_1=\kappa_1+\kappa_2, \qquad \chi_2=\kappa_1\kappa_2=1,
  \label{eq:G}
\end{equation}
matching \texttt{paper/xxxCG.tex} \S2.2 exactly, with numeric values reused verbatim from
\texttt{XXX\_CG\_L2\_Clean.wb}'s Cell 1: $\theta_1=1/3$, $\theta_2=1/7$, $h=1$,
$\kappa_1=e^{i\zeta(3)}$ (30-digit precision), $\kappa_2=1/\kappa_1=\overline{\kappa_1}$. The one new
parameter is the third-site inhomogeneity, $\theta_3=4/5$, chosen and mechanically confirmed in
Cell 2 (\texttt{IntegerQ[$\theta_3-\theta_1$]}, \texttt{IntegerQ[$\theta_3-\theta_2$]} both
\texttt{False}) so that no pair of inhomogeneities differs by an integer, avoiding an accidental
Baxter-equation degeneracy. Unlike L2, there is no dedicated numerical check cell for $G$ itself in
this notebook; its trace/determinant identity was already confirmed at L2 and the underlying twist
eigenvalues are reused unchanged.

\section{The three-site monodromy and transfer matrix}
\label{sec:twosite}

The three-site Lax operators (Cell 11, routines \texttt{L1}, \texttt{L2}, \texttt{L3}) act on
$V_{\lambda_1}\otimes V_{\lambda_2}\otimes V_{\lambda_3}$ by shifting each site's single-site Lax
matrix by its own inhomogeneity and tensoring with the identity on the other two factors,
\[
  L_1(u) = L^{[\lambda_1]}(u-\theta_1)\otimes \mathbbm{1}_{\lambda_2}\otimes\mathbbm{1}_{\lambda_3},
  \qquad\text{(similarly for $L_2$, $L_3$).}
\]
All three representation labels must be supplied to every routine, even though only one enters the
shifted Lax matrix directly, since the identity factors' sizes depend on the \emph{other two} sites'
representations; Cell 12 confirms $L_1(u)$ has dimension $8\times8$ for
$(\lambda_1,\lambda_2,\lambda_3)=(1,1,1)$ ($\dim\{1,0\}=2$ per site, $2^3=8$). The three-site
monodromy matrix (Cell 14, memoized routine \texttt{T}) is the ordered triple product in auxiliary
space,
\begin{equation}
  T(u) = L_1(u)\, L_2(u)\, L_3(u),
  \label{eq:monodromy}
\end{equation}
all three factors carrying the same spectral parameter $u$ (each site's shift already baked into
$L_1,L_2,L_3$); Cell 15 confirms the same $8\times8$ dimension.

The fundamental, companion-twisted transfer matrix (Cell 16),
\begin{equation}
  t_{1,1}(u) = \operatorname{tr}_a\!\bigl[G\,T(u)\bigr],
  \label{eq:t11}
\end{equation}
generalizes \texttt{paper/xxxCG.tex} \S2 to the three-site monodromy \eqref{eq:monodromy}, and is now
\textbf{cubic} in $u$ (three Lax factors), where its L2 counterpart was quadratic. Commutativity of
this family at different spectral parameters is checked in two steps, mirroring a pattern used
throughout Section A of this notebook: a symbolic check (Cell 17, over
$\lambda_1,\lambda_2,\lambda_3\in\{1,2\}$ --- a range trimmed from L2's $\{1,\dots,4\}$ to bound
compute cost, since the physical-space dimension grows as $(\lambda+1)^3$ rather than $(\lambda+1)^2$)
does \emph{not} literally collapse the commutator $t_{1,1}(u)t_{1,1}(v)-t_{1,1}(v)t_{1,1}(u)$ to the
singleton list $\{0\}$; the notebook's own Cell 18 comment attributes this to a benign
simplification artifact of mixing exact-rational polynomial coefficients with 30-digit-precision
numeric $\kappa_1,\kappa_2$ values, not an actual nonzero residual. The substantive confirmation is
Cell 18, which evaluates the same commutator at fixed generic complex spectral parameters
($u=1.234+0.567i$, $v=5.678-2.345i$, both carried at 30-digit precision) over the same $\lambda$
range and obtains exactly $0$ after \texttt{Chop}. This symbolic-then-numeric two-step pattern recurs
at the SoV covector check in \S\ref{sec:sov} below.

A purely symbolic diagonal-twist variant (Cell 19, routine \texttt{t3diag}), using formal indexed
eigenvalues $\kappa(1),\kappa(2)$ --- \textbf{distinct symbols} from the numeric plain twist
eigenvalues $\kappa_1,\kappa_2$ --- and the corresponding total (triple co-product) generator
(Cell 20, routine \texttt{EE}),
\[
  E_{ij}^{\mathrm{tot}} = E_{ij}\otimes\mathbbm{1}\otimes\mathbbm{1}
  + \mathbbm{1}\otimes E_{ij}\otimes\mathbbm{1}
  + \mathbbm{1}\otimes\mathbbm{1}\otimes E_{ij},
\]
are used only for a symbolic asymptotics cross-check, exactly as at L2. The central-charge identity
$E_{11}^{\mathrm{tot}}+E_{22}^{\mathrm{tot}}=(\lambda_1+\lambda_2+\lambda_3)\,\mathbbm{1}$ is checked
in Cell 21 over $\lambda_1,\lambda_2,\lambda_3\in\{1,2,3\}$ and collapses exactly to $\{0\}$.

\section{The load-bearing asymptotic cross-check}
\label{sec:asymptotics}

Cell 22 defines the predicted large-$u$ asymptotic expansion of the diagonal-twist transfer matrix up
to and including the $u^2$ term --- one order higher than L2's corresponding check, which only needed
the $u^1$ term, since $t_{1,1}(u)$ is now cubic rather than quadratic:
\begin{equation}
  t^{\mathrm{diag}}_{1,1}(u) \;\sim\; \bigl(\kappa(1)+\kappa(2)\bigr)u^3\,\mathbbm{1}
  - u^2\bigl(\kappa(1)+\kappa(2)\bigr)(\theta_1+\theta_2+\theta_3)\,\mathbbm{1}
  - u^2\,h\bigl(\kappa(1)\,E_{11}^{\mathrm{tot}} + \kappa(2)\,E_{22}^{\mathrm{tot}}\bigr).
  \label{eq:asympt}
\end{equation}
Cell 23 then confirms, for $(\lambda_1,\lambda_2,\lambda_3)=(1,1,1)$, that
\[
  \operatorname{Coefficient}\bigl[t^{\mathrm{diag}}_{1,1}(u) - \text{(RHS of \eqref{eq:asympt})},\, u,\, 2\bigr]
\]
collapses exactly to the $8\times8$ zero matrix. This is, by the notebook's own explicit comment on
this cell, \textbf{the single most important check in the whole notebook}: it establishes that the
$u^2$ coefficient $C_2$ of the genuine transfer matrix is a pure Cartan/total-generator quantity, and
therefore that the weight label $M$ of an eigenstate can be read off $C_2$'s eigenvalue algebraically
rather than treated as an unknown, free per-state number. Everything downstream in
\S\ref{sec:baxter} --- in particular \texttt{TauEigensystem}'s $M$-extraction formula --- depends on
this cross-check having passed; the notebook's comment is explicit that a nonzero result here would
mean the entire Baxter re-derivation needs to be revisited, not patched around downstream.

\section{Quantum determinant, fusion hierarchy, and the SoV covector basis}
\label{sec:sov}

The (untwisted) quantum determinant of the three-site monodromy (Cell 24, routine \texttt{qdetT}),
\begin{equation}
  \mathrm{qdet}\,T(u) = T_{11}(u)\,T_{22}(u-h) - T_{21}(u)\,T_{12}(u-h),
  \label{eq:qdet}
\end{equation}
and the Hirota (T-system) fusion hierarchy built from it (Cell 25) --- boundary conditions
$t_{0,s}\equiv t_{a,0}\equiv\mathbbm{1}$, the antisymmetric fusion
$t_{2,1}(u)=\chi_2\,\mathrm{qdet}\,T(u)$, and the symmetric-fusion Hirota recursion
$t_{1,s}(u)=t_{1,s-1}(u)\,t_{1,1}(u+(s-1)h)-t_{2,1}(u+(s-1)h)\,t_{1,s-2}(u)$ --- generalize
\texttt{paper/xxxCG.tex} \S2.2 to the three-site case with an algebraic form \emph{identical} to
L2's: only the domain (a triple $(\lambda_1,\lambda_2,\lambda_3)$, not a pair) changes. As at L2,
there is no standalone check cell for the fusion hierarchy; its correctness is implicit in the SoV
covector eigenvalue check below.

The $\qth$ factor for the three-site chain (Cell 26) is the triple product
$\qth(u)=(u-\theta_1)(u-\theta_2)(u-\theta_3)$, generalizing \texttt{paper/xxxCG.tex} \S2.3's
$\qth(u)=u-\theta$; unlike L2, this notebook drops the unused single-site factor $Q_{\theta,1}(u)$
that L2 kept purely for parity with the file it was ported from (the notebook's own comment on Cell
26 records this as a deliberate YAGNI simplification), and $\qth$ is defined only once in this
notebook --- Section B reuses it directly rather than redefining it, since this notebook was never
split into two separately self-contained source files the way L2's predecessors were.

The SoV $B$-operator for $L=3$ (Cell 28, routine \texttt{B3}) is again the bare $(1,1)$ monodromy
element,
\[
  B_3(u) = T_{11}(u),
\]
matching the single-site convention $B^{[\lambda]}(u)=L^{[\lambda]}_{11}(u-\theta)$ of
\texttt{paper/xxxCG.tex} \S2.3. The SoV pseudovacuum and covector basis (Cell 30) again tensors the
three single-site lowest-weight covectors,
\[
  \langle x_3(0,0,0)| = \langle x^{[\lambda_1]}_0|\otimes\langle x^{[\lambda_2]}_0|\otimes\langle x^{[\lambda_3]}_0|,
\]
with the full basis (memoized routine \texttt{x3}) obtained by acting with the fused transfer
matrices,
\begin{equation}
  \langle x_3(s_1,s_2,s_3)| =
  \frac{\langle x_3(0,0,0)|\;t_{1,s_1}(\theta_1)\;t_{1,s_2}(\theta_2)\;t_{1,s_3}(\theta_3)}
       {\displaystyle\prod_{k=1}^{s_1}\qth\bigl(\theta_1+(k-1-\lambda_1)h\bigr)\;
        \prod_{k=1}^{s_2}\qth\bigl(\theta_2+(k-1-\lambda_2)h\bigr)\;
        \prod_{k=1}^{s_3}\qth\bigl(\theta_3+(k-1-\lambda_3)h\bigr)},
  \label{eq:x3}
\end{equation}
generalizing \texttt{paper/xxxCG.tex} \S2.3 exactly the way L2's $x_2$ did, via independent fusion
actions and normalizations at each of the three sites. As with the commutativity check of
\S\ref{sec:twosite}, verification here again follows the symbolic-then-numeric two-step pattern:
Cell 31's symbolic \texttt{Simplify//Flatten//Union} (over $s_1\in\{0,\dots,\lambda_1\}$,
$s_2\in\{0,\dots,\lambda_2\}$, $s_3\in\{0,\dots,\lambda_3\}$, $\lambda_1,\lambda_2,\lambda_3\in\{1,2\}$,
again trimmed from L2's $\{1,2,3\}$ for compute cost) does not literally collapse the residual
\[
  \langle x_3(s_1,s_2,s_3)|\, B_3(u) - (u-\theta_1-hs_1)(u-\theta_2-hs_2)(u-\theta_3-hs_3)\,\langle x_3(s_1,s_2,s_3)|
\]
to $\{0\}$, for the same benign precision-tracked-arithmetic reason noted in Cell 32's own comment;
the substantive numeric confirmation is Cell 32, evaluating the same residual at a fixed generic
complex point ($u=3.7+1.2i$ at 30-digit precision) over the same range, designed to \texttt{Chop} down
to exactly $0$, confirming that $x_3$ diagonalizes $B_3$ with the product-of-three-per-site
eigenvalue expected from the tensor-product SoV structure.

Cell 32 originally shipped with a bracket-count bug (one stray closing bracket before
\texttt{//Chop}, making the cell a syntax error rather than a computation); after the fix was
applied its displayed output briefly went stale even though the corrected source had already been
saved. Both issues were caught and corrected in-session: the cell was re-run fresh and confirmed to
evaluate to exactly $0$, and that result is what is now saved on disk.

\section{The Baxter TQ equation}
\label{sec:baxter}

The routines in Section B of the notebook jointly solve the transfer-matrix eigenproblem
$t_{1,1}(u)\Psi=\tau(u)\Psi$ together with the Baxter TQ equation
\begin{equation}
  \kappa_1\,a(u)\,q(u+h) + \kappa_2\,\qth(u)\,q(u-h) = \tau(u)\,q(u),
  \label{eq:TQ}
\end{equation}
following the same two-stage strategy as L2 --- \texttt{TauEigensystem} diagonalizes the genuine
transfer matrix directly, then \texttt{QSolve} solves linearly for $q_n(u)$ at the now-known
eigenvalue data --- but with \emph{two} free per-state numbers $(J_1,J_0)$ in place of L2's single
$J$, as explained in the Overview and confirmed by the cross-check of \S\ref{sec:asymptotics}.
\texttt{paper/xxxCG.tex} does not yet state a Baxter/TQ equation, so none of this section has a paper
cross-reference.

\subsection*{Baxter-equation ingredients}

Cell 34 defines the Bethe/Baxter coefficient
\[
  a(u) = (u-\theta_1-h\lambda_1)(u-\theta_2-h\lambda_2)(u-\theta_3-h\lambda_3),
\]
reuses $\qth(u)=(u-\theta_1)(u-\theta_2)(u-\theta_3)$ from \S\ref{sec:sov} without redefinition, and
the state-independent part of the transfer-matrix eigenvalue,
\[
  \tau_0(u) = (\kappa_1+\kappa_2)u^3 - (\theta_1+\theta_2+\theta_3)(\kappa_1+\kappa_2)u^2
    - u^2\,h\bigl(\kappa_1(\lambda_1+\lambda_2+\lambda_3-M)+\kappa_2 M\bigr),
\]
which by construction has only $u^3,u^2$ terms --- zero $u^1,u^0$ coefficients --- so that the full
eigenvalue is $\tau=\tau_0+J_1u+J_0$ with $(J_1,J_0)$ the two free per-state numbers introduced above.
The helper \texttt{qpoly[clist,uu]} evaluates a polynomial from an ascending coefficient list, formed
in a dummy symbol first to avoid an indeterminate $0^0$ at $uu=0$; it is unchanged from L2's. This is
the standard XXX Baxter-equation form generalized to three inhomogeneous sites with a companion twist
of eigenvalues $\kappa_1,\kappa_2$.

\subsection*{The Baxter matrix and transfer-matrix coefficients}

\texttt{BaxterMatrix} at fixed $(\lambda_1,\lambda_2,\lambda_3,M)$ (Cell 35) builds, for
$k=0,\dots,M$, the column obtained by expanding
\begin{equation}
  O[u^k] = \kappa_1\,a(u)(u+h)^k + \kappa_2\,\qth(u)(u-h)^k - \tau_0(u)\,u^k
  \label{eq:baxtercol}
\end{equation}
as a polynomial in $u$ (now degree $k+3$, one degree higher than L2's, since $a,\qth,\tau_0$ are
cubic here rather than quadratic); the routine hard-asserts that the degree-$(M{+}2)$ and
degree-$(M{+}3)$ coefficients vanish for every column, and the surviving degree-$0,\dots,M$
coefficients become the columns of the resulting $(M{+}1)\times(M{+}1)$ matrix. Cell 35's smoke test,
\texttt{Dimensions /@ (BaxterMatrix[1,1,1,\#]\&/@\{0,1,2,3\})}, gives $\{\{1,1\},\{2,2\},\{3,3\},\{4,4\}\}$
with no assertion failure.

The choice of \emph{which} two high-degree coefficients to assert deserves its own explanation, since
it is the one place where L3's structure genuinely differs from a mechanical shift of L2's. Writing
$O_0[u^k]$ for the $\tau_0$-only construction \eqref{eq:baxtercol} used by \texttt{BaxterMatrix}, and
$O_{\mathrm{full}}[u^k]$ for the corresponding expression built from the true eigenvalue
$\tau=\tau_0+J_1u+J_0$, the two differ by
\[
  O_0[u^k] = O_{\mathrm{full}}[u^k] + J_1\,u^{k+1} + J_0\,u^k.
\]
The $J_0\,u^k$ term contaminates row $=$ column (the diagonal, rows $0,\dots,M$), which
\texttt{QSolve}'s $J_0\cdot\mathbbm{1}$ subtraction undoes later. The $J_1\,u^{k+1}$ term contaminates
row $k{+}1$ relative to column $k$; for the top column $k=M$ this lands exactly at row $M{+}1$ ---
which is therefore \emph{not} a pure per-column identity at $L=3$, unlike $L=2$ (which has no
$J_1$-like free parameter feeding an off-diagonal contamination at all). Row $M{+}1$'s physical
content is supplied later by \texttt{QSolve}'s $J_1\cdot(\text{shiftDown})$ correction, and is
validated only by the end-to-end residual check of \S\ref{sec:endtoend}, not by a symbolic assert
here. Rows $M{+}2$ and $M{+}3$ are never reached by either contamination and remain clean per-column
identities --- hence exactly two hard-asserted rows, the same count as L2, just shifted up one degree.
This generalizes cleanly to any chain length $L$: exactly the top two extra-degree coefficients,
$M{+}(L{-}1)$ and $M{+}L$, are always pure per-column identities regardless of $L$, while any lower
extra-degree row is contaminated by one of the $L{-}1$ free state-dependent parameters and belongs to
\texttt{QSolve} instead.

Separately, \texttt{t3Coeffs[$\lambda_1$,$\lambda_2$,$\lambda_3$]} (Cell 36, memoized) extracts the
coefficient matrices $C_0,C_1,C_2,C_3$ of $t_{1,1}(u)=C_0+C_1u+C_2u^2+C_3u^3$ entrywise via
\texttt{CoefficientList}, returning a four-element list (L2's \texttt{t2Coeffs} returned three). Its
own check confirms $C_3=(\kappa_1+\kappa_2)\mathbbm{1}$ (the transfer matrix's leading-order,
twist-independent asymptotics), obtaining $\max|C_3-(\kappa_1+\kappa_2)\mathbbm{1}|\to0\times10^{-29}$
for $(\lambda_1,\lambda_2,\lambda_3)=(1,1,1)$.

The predicted dimension of the weight-$M$ subspace of $V_{\lambda_1}\otimes V_{\lambda_2}\otimes
V_{\lambda_3}$ (Cell 37, routine \texttt{dPred}) is now a genuine triple convolution --- the number of
$(m_1,m_2,m_3)$ with $m_1+m_2+m_3=M$, $0\le m_i\le\lambda_i$ --- rather than L2's closed form
$\min(M,\lambda_1,\lambda_2,\lambda_1+\lambda_2-M)+1$; there is no closed form at $L=3$, and the
notebook computes it by a direct double sum. This has a concrete downstream consequence: L2's
\texttt{dPred} is always $\le M+1$ by construction ($M$ is one of its own \texttt{Min} arguments), but
L3's \texttt{dPred} can \emph{exceed} $M+1$ --- for instance $d_{\mathrm{pred}}(1,1,1,1)=3>M+1=2$, a
binomial-like distribution of states across three sites. Consequently, every accessor below
(\texttt{Psi}, $\tau$, \texttt{QSolve}) deliberately omits the redundant-at-L2 gross bound $n\le M$,
relying solely on the explicit multiplicity check $n>d_{\mathrm{pred}}-1$ instead.

\subsection*{Diagonalizing the genuine transfer matrix}
\label{sec:taueigensystem}

\texttt{TauEigensystem[$\lambda_1$,$\lambda_2$,$\lambda_3$]} (Cell 38, memoized) diagonalizes the
genuine transfer matrix directly. Since $C_0,C_1,C_2,C_3$ pairwise commute and $C_3$ is already known
to be scalar, the routine diagonalizes a generic \emph{two}-parameter combination
$C_0+r_1C_1+r_2C_2$ (L2 needed only one free parameter $r$), trying six candidate pairs
$(r_1,r_2)\in\{(7/3,11/5),(\pi,e),\dots\}$ until the combination's eigenvectors have full rank
$d=(\lambda_1+1)(\lambda_2+1)(\lambda_3+1)$ with $d$ distinct eigenvalues. For each eigenvector $v$ it
extracts
\[
  t_0=\frac{v\cdot C_0v}{v\cdot v}, \qquad t_1=\frac{v\cdot C_1v}{v\cdot v}, \qquad t_2=\frac{v\cdot C_2v}{v\cdot v},
\]
and hard-aborts if any of the three residuals $\|C_iv-t_iv\|$ ($i=0,1,2$) exceeds $10^{-15}$. The
weight label $M$ is read off algebraically from $t_2$ --- exactly as L2 read $M$ from $t_1$, one order
down since the polynomial is one degree higher --- by inverting $\tau_0$'s $u^2$ coefficient,
\begin{equation}
  M = \frac{-(\theta_1+\theta_2+\theta_3)(\kappa_1+\kappa_2) - h\kappa_1(\lambda_1+\lambda_2+\lambda_3) - t_2}
           {h(\kappa_2-\kappa_1)},
  \label{eq:Mextract}
\end{equation}
hard-asserted to lie within $10^{-10}$ of an integer. Unlike $M$, the remaining two numbers
$(J_1,J_0)=(t_1,t_0)$ are recorded \emph{directly per eigenvector}, with no closed-form formula ---
the same "read it straight off the eigenvector" treatment L2 gave its single $J=t_0$, now applied to
a pair. States within each $M$-sector are sorted by $(\operatorname{Re}J_0,\operatorname{Im}J_0,
\operatorname{Re}J_1,\operatorname{Im}J_1)$ to assign the index $n$, and each sector's count is
hard-asserted to equal $d_{\mathrm{pred}}(M,\lambda_1,\lambda_2,\lambda_3)$ exactly. The result is an
association keyed by $(M,n)$ giving $(J_{1,n},J_{0,n},\Psi_{M,n})$, with $\Psi_{M,n}$ normalized so
its last component is $1$.

These in-function hard asserts (rank/degeneracy resolution over the two-parameter combination,
residual below $10^{-15}$ against all three of $C_0,C_1,C_2$, integer $M$ within $10^{-10}$, exact
per-sector counts) constitute the check; the smoke test in Cell 38 evaluates
\texttt{KeySort[TauEigensystem[1,1,1]][[All,"J0"]]} for all $8$ states with no error, e.g.
\begin{verbatim}
KeySort[TauEigensystem[1, 1, 1]][[All, "J0"]] // N
(* <|{0,0}->-1.0023661289808585-2.5229544323858897 I,
     {1,0}->-1.9591204341670663-0.8728625487854049 I,
     {1,1}->-0.08422543749008554+0.27639976349988676 I,
     {1,2}->1.0375469819606056-0.9937091985069969 I,
     {2,0}->-1.9591204341670663+0.8728625487854049 I,
     {2,1}->-0.08422543749008554-0.27639976349988676 I,
     {2,2}->1.0375469819606056+0.9937091985069969 I,
     {3,0}->-1.0023661289808585+2.5229544323858897 I|> *)
\end{verbatim}
Note that the $M=1$ and $M=2$ sectors each carry \emph{three} states, $d_{\mathrm{pred}}(1,1,1,1)=3$
exceeding $M+1=2$ --- a concrete illustration of the \texttt{dPred} behavior noted above.
Independently, the end-to-end TQ residual check of \S\ref{sec:endtoend} re-confirms
\texttt{TauEigensystem}'s output from a completely different derivation path (through
\texttt{QSolve}/\texttt{BaxterMatrix} rather than this $C_0,C_1,C_2$ extraction).

Bounds-checked accessors \texttt{Psi[$\lambda_1$,$\lambda_2$,$\lambda_3$][M,n]} and
$\tau[\lambda_1,\lambda_2,\lambda_3][M,n][u]=\tau_0(u)+J_{1,n}u+J_{0,n}$ (Cell 39) read the shared
cache; both require $M$ an integer in $0,\dots,\lambda_1+\lambda_2+\lambda_3$ and $n$ a non-negative
integer (else a bounds message), and, if $n$ is in this gross range but exceeds
$d_{\mathrm{pred}}-1$, fire a spurious-state message and return \texttt{Missing[...]}. A deliberate
smoke test calls \texttt{Psi[2,1,1][3,3]} --- $d_{\mathrm{pred}}(3,2,1,1)=3$, so only $n=0,1,2$ exist
even though $n=3\le M=3$ passes the gross bound --- and confirms the expected spurious-state message
and missing result.

\subsection*{Solving for the Q-function}

\texttt{QSolve[$\lambda_1$,$\lambda_2$,$\lambda_3$][M,n]} (Cell 40, memoized) solves \eqref{eq:TQ} as
a \emph{linear} problem, since both $J_1$ and $J_0$ (hence $\tau$) are already known from
\texttt{TauEigensystem}: it finds the monic degree-$M$ polynomial $q_n$ satisfying
\[
  \Bigl(\mathrm{BaxterMatrix}(\lambda_1,\lambda_2,\lambda_3,M) - J_{1,n}\,\Sigma - J_{0,n}\,\mathbbm{1}\Bigr)\, q_n = 0,
\]
where $\Sigma$ is the $(M{+}1)\times(M{+}1)$ subdiagonal shift matrix (row $=$ column $+1$),
representing multiplication of the underlying $u^k$ monomial by an extra factor of $u$ --- i.e.\ the
term $J_1u^{k+1}$ lands one row below the diagonal relative to column $k$. L2's \texttt{QSolve} only
needed a single $J\cdot\mathbbm{1}$ subtraction, since its $\tau$ carried only the one free constant
term. The same precision fix as L2 is used before computing the nullspace: the matrix is
\texttt{Chop}'d at threshold $10^{-12}$ (converting precision-tracked-but-exact zero entries, which a
bare nullspace call would otherwise treat as unknown, into confirmed zeros) with an explicit
tolerance of $10^{-10}$ kept in addition. As a representative example, reproduced from Cell 40,
\begin{verbatim}
QSolve[1, 1, 1][2, 0]
(* {-0.04995738064123986793660273-1.230053239147504224150097 I,
    -1.856051357929965424850272+1.317538826427866191621732 I,
    1.0000000000000000000000000+0.×10^-26 I} *)
\end{verbatim}
gives the ascending coefficients of the monic degree-2 Q-polynomial for that sector. Unlike L2, this
notebook has no standalone regression-guard cell isolating a historically-critical precision case;
end-to-end correctness of $q_n$ across every $(M,n)$ is confirmed instead by the residual check of
\S\ref{sec:endtoend}.

Finally, $Q_{M,n}(u)=\kappa_1^{u/h}\,q_n(u)$ (routine \texttt{Q1}, Cell 41) assembles the actual
Baxter Q-function value from the same cache, mutually consistent with $\tau_{M,n}$ by construction. A
deliberate spurious-index smoke test in Cell 42, \texttt{Q1[1,1,1][0,1][u]}
($d_{\mathrm{pred}}(0,1,1,1)=1$, so $n=1$ does not exist at $M=0$), resolves correctly to
\texttt{\$Failed} internally (confirmed by this agent via direct kernel-state inspection) though the
cell's own rendered output displays nothing visible rather than an explicit failure message or value
--- a cosmetic tool-capture nuance discussed further in \S\ref{sec:gotchas}, not a correctness issue.
A valid example, from Cell 41,
\begin{verbatim}
{Q1[1, 1, 1][1, 0][2], N[\[Tau][1, 1, 1][1, 0][2], 10]}
(* {-1.555304313657261488454067-0.106068184694823126416499 I,
    2.455604872-1.170833864 I} *)
\end{verbatim}

\subsection*{End-to-end TQ residual regression check}
\label{sec:endtoend}

The notebook closes with its own complete, first-class validation of the Baxter TQ equation
\eqref{eq:TQ}, assembled purely from \texttt{Q1} and $\tau$ --- two quantities tracing back to
genuinely different derivation paths (\texttt{Q1} through \texttt{QSolve}/\texttt{BaxterMatrix},
$\tau$ through \texttt{TauEigensystem}'s $C_0,C_1,C_2$ diagonalization), so a small residual is strong
evidence that both derivations, including the re-derived $L=3$ $M$-extraction and $(J_1,J_0)$
structure, are correct. At $L=2$ this check exists only in the separate
\texttt{Experiments/Baxter\_Solver\_L2.wb}; here it is part of the Clean notebook itself, per this
notebook's own "no dropped validation" design constraint. The residual function (Cell 43) is
\[
  \mathrm{checkTQ}(\lambda_1,\lambda_2,\lambda_3,M,n) = \max_{u\in\{0,1,2,-1,5\}}
  \Bigl|\,a(u)\,Q(u+h) - \tau(u)\,Q(u) + \qth(u)\,Q(u-h)\,\Bigr|,
\]
sampled at points avoiding the roots of $a$ and $\qth$. The outer sampling loop deliberately uses a
distinct dummy variable rather than the bare global symbol \texttt{u} that \texttt{BaxterMatrix} uses
internally for its own \texttt{CoefficientList} call (not \texttt{Module}-localized there) --- using
the same symbol in both places risks \texttt{Table}'s dynamic variable binding clobbering
\texttt{BaxterMatrix}'s internal symbol during an uncached nested call, a bug the L2 predecessor
already avoided the same way.

Cell 43 checks every $(M,n)$ at $(\lambda_1,\lambda_2,\lambda_3)=(1,1,1)$: all 8 residuals land
between $\sim10^{-23}$ and $\sim10^{-19}$. Cell 44 checks the asymmetric case
$(\lambda_1,\lambda_2,\lambda_3)=(2,1,1)$: all 12 residuals land between $\sim10^{-23}$ and
$\sim10^{-18}$. Both are exact zero at the tracked $\sim30$-digit numeric precision budget:
\begin{verbatim}
Table[{M, n, checkTQ[1, 1, 1, M, n]}, {M, 0, 3}, {n, 0, dPred[M, 1, 1, 1] - 1}] // Flatten[#, 1] &
(* {{0,0,0.×10^-23}, {1,0,0.×10^-22}, {1,1,0.×10^-22}, {1,2,0.×10^-22},
    {2,0,0.×10^-21}, {2,1,0.×10^-21}, {2,2,0.×10^-21}, {3,0,0.×10^-19}} *)
\end{verbatim}
This matches exactly what the task that produced this notebook set out to confirm: the end-to-end TQ
residual is exact zero at 18--23 tracked digits for both $(\lambda_1,\lambda_2,\lambda_3)=(1,1,1)$ and
$(2,1,1)$.

\section{Known gotchas and non-obvious conventions}
\label{sec:gotchas}

\begin{itemize}
\item \textbf{Cell 32 originally shipped with a bracket-count bug (fixed).} The numeric confirmation
  that $x_3$ diagonalizes $B_3$ at a fixed generic complex point (closing Section A) originally had
  one stray extra closing bracket before \texttt{//Chop}, making the cell a syntax error rather than
  a computation; its displayed output then briefly went stale after the fix was applied. Both were
  caught and corrected in-session --- the cell was re-run fresh and confirmed to evaluate to exactly
  $0$, and that result is what is saved on disk now.
\item \textbf{\texttt{Q1[1,1,1][0,1][u]} (Cell 42) displays no visible output.} This deliberate
  spurious-index smoke test resolves correctly to \texttt{\$Failed} internally (confirmed via direct
  kernel-state inspection), but its own rendered cell output shows nothing displayed rather than an
  explicit value or message text. This looks like a cosmetic tool-capture/rendering nuance,
  comparable to the display glitch recorded for \texttt{XXX\_CG\_L2\_Clean.wb}, not a correctness
  problem.
\item \textbf{Symbolic \texttt{Simplify} checks don't literally collapse to $\{0\}$ at $L=3$.} Both
  the transfer-matrix commutativity check (\S\ref{sec:twosite}) and the SoV covector eigenvalue check
  (\S\ref{sec:sov}) fail to simplify their symbolic residual down to the literal singleton $\{0\}$,
  unlike their L2 counterparts. The notebook's own comments attribute this to mixed exact-rational
  and 30-digit-precision numeric arithmetic that \texttt{Simplify} does not always fully merge, not
  to an actual nonzero residual; each symbolic check is followed by a substantive numeric
  confirmation at a fixed generic complex point that does \texttt{Chop} down to exactly $0$.
\item \textbf{\texttt{dPred} can exceed $M+1$ -- no gross $n\le M$ bound anywhere in Section B.}
  Unlike L2 (where \texttt{dPred} is always $\le M+1$ by construction), L3's \texttt{dPred} is a
  genuine triple convolution that can exceed $M+1$ (e.g.\ $d_{\mathrm{pred}}(1,1,1,1)=3>2$).
  \texttt{Psi}, $\tau$, and \texttt{QSolve} therefore gross-check only $n\ge0$, relying entirely on
  the explicit multiplicity check as the sole source of truth for which $n$ actually exist per
  sector.
\item \textbf{\texttt{BaxterMatrix}'s hard assert checks degrees $M{+}2,M{+}3$, not $M{+}1$.} Degree
  $M{+}1$ is genuinely state-dependent (contaminated by $J_1$) at $L=3$ and must not be asserted to
  vanish inside \texttt{BaxterMatrix} -- its correct handling is \texttt{QSolve}'s job via the
  subdiagonal shift correction, validated only by the end-to-end residual check
  (\S\ref{sec:endtoend}), not by a symbolic assert. This generalizes to any chain length $L$: exactly
  the top two extra-degree coefficients are always pure identities, regardless of $L$.
\item \textbf{\texttt{QSolve} subtracts two terms, not one.} $J_{1,n}\Sigma + J_{0,n}\mathbbm{1}$,
  both subtracted from \texttt{BaxterMatrix} before the nullspace solve -- L2's \texttt{QSolve} only
  needed $J\cdot\mathbbm{1}$. The chop-before-nullspace precision fix is otherwise unchanged from L2.
\item \textbf{Memoization caches stale results.} Several routines throughout the notebook
  (\texttt{ps}, \texttt{dim}, the GT generator routines, \texttt{T}, the fused \texttt{t3[1,s]},
  \texttt{x3}, \texttt{t3Coeffs}, \texttt{TauEigensystem}, \texttt{QSolve}) are self-memoizing. After
  fixing any upstream definition, re-run the defining cell and everything downstream that used the
  stale cache.
\item \textbf{High-precision twist reused verbatim from L2; one new parameter.} $h=1$ and
  $\kappa_1=e^{i\zeta(3)}$ at 30-digit precision (with $\kappa_2=1/\kappa_1=\overline{\kappa_1}$) are
  reused unchanged from L2, as are $\theta_1=1/3,\theta_2=1/7$; only $\theta_3=4/5$ is new, chosen
  and mechanically confirmed (Cell 2) to differ from $\theta_1,\theta_2$ by a non-integer amount,
  avoiding an accidental Baxter-equation degeneracy. Numeric checks throughout rely on this tracked
  precision -- residuals reported at $10^{-18}$ to $10^{-29}$ are relative to this roughly 30-digit
  budget, not machine precision.
\item \textbf{Compute-cost-driven range trimming vs.\ L2.} Several Section A validation loops use a
  smaller $\lambda$-range than their L2 counterparts (e.g.\ $\lambda_i\in\{1,2\}$ where L2 used
  $\{1,\dots,4\}$ or $\{1,2,3\}$), because the physical-space dimension grows as $(\lambda+1)^3$ at
  three sites rather than $(\lambda+1)^2$ at two -- a deliberate, commented compute-cost tradeoff, not
  a reduction in coverage (both trimmed-range endpoints still exercise the general, unequal-$\lambda$
  case).
\end{itemize}

\end{document}
